\documentclass{scrbook}
\usepackage[utf8]{inputenc}
\usepackage[T1]{fontenc}
\usepackage{geometry}
\usepackage{graphicx}
\usepackage{float}
\usepackage{hyperref}
\usepackage[english,indonesian]{babel}

\title{Optimasi Pemrosesan Template LaTeX Menggunakan Pendekatan Berbasis AI}
\author{Academic Support Agent}
\date{\today}

\begin{document}
\maketitle

\begin{abstract}

\end{abstract}

\chapter{PENDAHULUAN}


\chapter{TINJAUAN PUSTAKA}


% Original content (commented out):
% % Clase de documento (Utilizar la que interese)
\documentclass{scrbook}

% MÁRGENES DE LAS PÁGINAS
\usepackage[
  inner	=	3.0cm, % Margen interior
  outer	=	2.5cm, % Margen exterior
  top	=	2.5cm, % Margen superior
  bottom=	2.5cm, % Margen inferior
  includeheadfoot, % Incluye cabecera y pie de página en los márgenes
]{geometry}
% Valor de interlineado
\renewcommand{\baselinestretch}{1.0} % 1 línea de interlineado

% Paquetes imprescindibles para la portada
\usepackage{fontspec} % Para poder importar las fuentes
\usepackage{textpos} % Para posicionar bloques de texto
\usepackage{graphicx} % Para cargar imágenes
\usepackage{setspace} % Para modificar los espacios arriba y abajo


% Universidad y facultad
\newcommand{\Universidad}{Universitat Politècnica de València}
\newcommand{\LogoUniversidad}{logos-universidad/LogoUPV}
\newcommand{\Facultad}{Escola Tècnica Superior d'Enginyeria Industrial}
\newcommand{\LogoFacultad}{logos-universidad/Logo_ETSII}
% Titulación
\newcommand{\Titulacion}{Grau en Enginyeria en Tecnologies Industrials}
% Título del trabajo
\newcommand{\titulo}{Título de la Prátctica}
% Tipo de trabajo
\newcommand{\tipotrabajo}{Informe de Prácticas}
% Datos del autor
\newcommand{\miNombre}{Nombre Apellido1 Apellido2}
% Datos del tutor/es
\newcommand{\miTutor}{Nombre Apellido1 Apellido2}
% Ubicación
\newcommand{\miUbicacion}{València}


\begin{document}


%%%%% PORTADA
\input{portada_bn} % Portada B/N

%%%%% RESUMEN
\input{resumen}

%%%%% INDICES
\tableofcontents

%%%%% Resto del documento
\chapter{Esto es un capítulo}

\section{Esto es una sección}

\subsection{Esto es una subsección}

\subsubsection{Esto es una subsubsección}

\end{document}


\end{document}